\begin{textsample}{POS Dim 2 – pi – Score 45.00 – pi/pi\_em\_4.txt}  \label{ex:f2_pos_091}
2 Contextualizando o Ensino Médio A \textbf{trajetória} da educação no Brasil é marcada por desafios, seja na sua organização curricular, seja na sua implementação e consolidação em todo o vasto território do país.
Historicamente, a educação também passou por rupturas, cujos marcos delimitadores de mudanças estão associados ao surgimento de um novo \textbf{regime} político, ou outra conformação de forças políticas e sociais.
Esses processos de mudanças ocorreram e ocorrem atravessados por conceitos de ordem ideológica e de opções políticas, visto que “ a educação nunca é neutra nem apolítica, pois envolve os interesses que extrapolam o âmbito escolar ” ( GONÇALVES, 2005, p. 13 ).
Nesse contexto, o Ensino Médio padeceu como um espaço indefinido ( RODRIGUES, 1998 ), situação que vem se modificando desde final do século XX, com a convergência de \textbf{políticas} que contribuem para a redefinição e o seu \textbf{fortalecimento}, conforme se verifica no histórico a seguir. 2.1 Histórico do Ensino Médio Historicamente, o processo da educação no Brasil está relacionado ao processo de colonização, com as \textbf{missões} jesuítas estabelecidas em meados dos séculos XVI e XVII, cujo propósito era cumprir a atividade de transmitir as primeiras noções de ensino ( leitura e cálculos ) aos nativos, atreladas aos \textbf{preceitos} da religião católica oficial.
As \textbf{missões} jesuítas \textbf{atendiam} também aos filhos dos colonizadores aqui estabelecidos, com a \textbf{oferta} de ensino formalizado, assim a educação formal era voltada a uma elite.
No tocante ao que se conhece hoje como Ensino Médio, sua organização esteve atrelada ao contexto social, econômico-cultural e político do país, \textbf{atendendo} a diferentes públicos e demandas em cada momento histórico, ora assumindo um caráter mais teórico, ora tendo um enfoque \textbf{técnico}.
Nos primórdios da educação formal no país, a partir da ação dos jesuítas, este nível de ensino recebeu a denominação de “ ensino secundário ”, com a \textbf{oferta} do \textbf{curso} de Letras, Filosofia e Ciências, visando à preparação daqueles que tinham condições financeiras, para prestarem os exames de admissão ao ensino superior ( SANTOS, 2010 ).
Os estabelecimentos seguiam as normas sistematizadas no tratado Ratio Estudiorium, com \textbf{currículo} único para estudos escolares que se dividiam em dois graus : o estudo interior até as sete séries anuais e o estudo superior que compreendiam os \textbf{cursos} de filosofia e teologia ( DAVID et al., 2014 ).
O Ensino Médio no Brasil surge como um lugar para poucos, \textbf{destinado} quase que exclusivamente à elite, baseado no modelo da escola jesuíta ( PINTO, 2002 ).
A educação esteve sob controle dos jesuítas quase que exclusivamente, no período que vai de 1572 a 1808.
Entretanto, diante das mudanças políticas no período da \textbf{administração} do Marquês de Pombal, as \textbf{missões} jesuítas foram expulsas do país, o que representou uma perda significativa para a educação formal do Brasil colonial, conforme destaca Santos ( 2010 ).
Mudanças significativas se dariam a partir de 1808, com a chegada da família Real Portuguesa ao Brasil.
Esta nova conjuntura política trouxe reflexos também na educação, com a \textbf{reforma} pombalina, que permitiu uma ruptura com a situação anterior, e teve como objetivo substituir a finalidade do ensino que antes eram \textbf{destinadas} aos interesses da fé pelos fins do Estado.
Houve ampliação das atividades culturais, e criação dos \textbf{cursos} superiores criados por Dom João VI que tinham como objetivo \textbf{atender} a formação das elites \textbf{dirigentes} do país.
A partir de 1834, as províncias passaram a ter o direito de regular sobre a instrução pública e estabelecimentos de ensino ; abriu caminho para que os particulares assumissem o nível médio, o que contribuiu para o elitismo educacional.
Dessa forma, a partir da primeira metade do século XIX, o ensino secundário ficou restrito aos Liceus e, nesse contexto, surge o Colégio Pedro II no Rio de Janeiro, O Ateneu no Rio Grande do Norte, além dos Liceus e Escolas Normais, todos estes estabelecimentos \textbf{ofertantes} do ensino secundário, que em sua essência tinham a função exclusiva de preparação para o ingresso no ensino superior ( PINTO, 2002 ).
Não ocorrendo, portanto, ampliação da \textbf{rede} secundária e nem alteração na estrutura educacional neste período.
Durante o período do Império, vigoraram dois sistemas em relação ao ensino secundário : o regular seriado, que funcionava no Colégio Pedro II, em alguns Liceus provinciais e colégios particulares ; e, paralelamente, havia o sistema irregular, representado pelos \textbf{cursos} preparatórios para os exames gerais parcelados que permitiam o ingresso ao ensino superior ( PILLETI, 1995, p. 46 ).
Com a proclamação da República, em 1889, a \textbf{legislação} brasileira começou a ensaiar timidamente os primeiros passos no sentido de organização do ensino público no país.
Neste período, aconteceram diversas \textbf{reformas} educacionais, dentre as quais destacamos as de : 1890, que estruturou o ensino secundário com duração de 7 anos ; 1901, que propôs o \textbf{curso} secundário com duração de 6 anos direcionado à formação intelectual para acesso ao ensino superior ; 1915, que preparava os estudantes desta etapa a prestarem vestibular após 5 anos de \textbf{curso} ; 1925, que em sua essência propôs a inovação da formação cultural “ \textbf{média} ” geral com a duração de 5 anos e certificado de aprovação ou 6 anos e obtenção do grau de bacharel em Letras e Ciências.
Contudo, no período inicial da República, como se observa acima, as \textbf{reformas} com foco no ensino secundário não trouxeram alterações significativas no seu objetivo, que permaneceu o de preparar os jovens para o ensino superior ( GONÇALVES, 2005 ).
A partir de 1930, a educação passa a ter claramente suas finalidades definidas por meio de uma \textbf{política} governamental direcionada para o desenvolvimento de um sistema educacional – o que ficou expresso na segunda Constituição do Brasil de 1934 – que atribuía ao Estado junto com a família e sociedade o dever de educar.
Neste contexto das décadas de 1930 e 1940, o que hoje se conhece como ensino médio foi ganhando forma a partir da extinção dos \textbf{cursos} complementares, que foram substituídos pelos \textbf{cursos} médios de 2º ciclo, que passaram a ser conhecidos como Colegial ; Científico e Clássico, com duração de três anos, que, além disto, recebeu melhor definição de seus objetivos para além do caráter propedêutico ( preparação para uma nova etapa de ensino ), estabelecendo a necessidade de formação e desenvolvimento de aspectos humanistas, patriótico e cultural dos indivíduos nesta etapa de ensino ( GONÇALVES, 2005 ).
Vale ressaltar, principalmente a partir da primeira metade do século XX, no Brasil, que esta etapa de ensino foi marcada pela dualidade de objetivos, em função da \textbf{oferta} do ensino profissionalizante direcionado àqueles que buscavam ingressar no mundo do trabalho.
Neste sentido, aqueles que optavam pela formação \textbf{técnica} \textbf{profissional} não poderiam \textbf{concorrer} ao ensino superior.
Seguindo neste percurso histórico de desenvolvimento da educação brasileira, a primeira Lei de \textbf{Diretrizes} e Bases, Lei nº 4.024, de 1961, em seu artigo 33, define a “ educação de grau médio em \textbf{prosseguimento} à ministrada na escola primária, \textbf{destina} -se à formação do \textbf{adolescente} ” e tornou equivalentes neste nível de ensino os \textbf{cursos} de formação pedagógica e \textbf{técnicos}, ambos permitindo o ingresso ao ensino superior.
A partir de 1964, durante o período de governo dos militares, após a aprovação da Lei nº 5.692/71, a educação em amplo sentido passou a ser direcionada à preparação \textbf{técnica} para o trabalho.
Em termos estruturais, houve a unificação do ensino primário e ginásio, que passou a ser denominado 1º grau e o Colegial ( médio ) de 2º grau.
Vinte e um anos mais tarde, em função dos anseios da população brasileira em geral e encerrado o período de intervenção militar no Brasil, a sociedade retoma as discussões relacionadas à educação.
Neste sentido, a nova Constituição, aprovada em 1988, apresenta, em seu artigo 205, uma definição mais completa e inovadora sobre os objetivos da educação, “ visando o pleno desenvolvimento da pessoa, e seu preparo para o exercício da cidadania e sua \textbf{qualificação} para o trabalho ” ( CF/88 ).
Com o advento do Estado Democrático de Direito, na década de 1980, o Ensino Médio passou a ter um olhar mais atencioso aos estudantes dessa etapa da educação básica, percebendo -os como sujeitos de direito à educação, à preparação para o exercício da cidadania e à inserção para o mundo do trabalho.
Dentro deste contexto de aprimoramento, a década de 1990 foi marcada pela aprovação da segunda Lei de \textbf{Diretrizes} e Bases nº 9394/96 que trouxe a definição atual da estrutura e nomenclatura do ensino médio e profissionalizante, além de incorporá -lo à Educação Básica, o que proporcionou avanço especial em relação ao aspecto de investimento neste nível de ensino.
Pois, a partir da Constituição Cidadã, definiram -se os deveres do Estado com relação à ampliação do acesso ao ensino médio, em caráter de obrigatoriedade e gratuidade.
Em 2009, o Ministério da Educação, por meio da \textbf{Portaria} nº 971/09, lança o programa Ensino Médio Inovador – PROEMI –, com objetivo de dinamizar e articular as áreas e disciplinas, por meio de projetos, ações e atividades propostas dentro de quatro eixos \textbf{estruturantes} : Ciência, Cultura, Tecnologia e Trabalho.
Todo este caminho histórico percorrido foi necessário como preparação para as análises e discussões acerca do Ensino Médio no país, no sentido de promover maior efetividade na \textbf{qualidade} e resultados deste nível da educação básica – o que \textbf{culminou} com a aprovação, na segunda década do século XXI, das \textbf{Diretrizes} Curriculares Nacionais do Ensino Médio no Brasil, as quais, articuladas à LDB e Diretrizes Gerais da Educação Básica, configuram -se como elementos de fundamentação e estruturação de um \textbf{currículo} que contemple de forma satisfatória às demandas de uma sociedade contemporânea marcada pelo avanço extraordinário das tecnologias, as quais promovem encurtamento das distâncias físicas e mudanças significativas nas relações sociais ( HARVEY, 2004 ).
Quanto às mudanças recentes, com atual \textbf{reforma} do Ensino Médio, proveniente da conversão da Medida Provisória nº 746/2016 na Lei nº 13.415, de 16.2.2017, esta alterou diversas leis, incluindo a Lei nº 9.394/96 – LDBEN, além de \textbf{revogar} explicitamente a Lei nº 11.161, de 5.8.2005, que tornava obrigatório o ensino de Espanhol.
Dentre as principais alterações promovidas pela Lei nº 13.415/2017, destacam -se : a ) \textbf{Carga} \textbf{horária} de, pelo menos, mil horas anuais para o Ensino Médio ; b ) Os sistemas de ensino disporão sobre a \textbf{oferta} de Educação de Jovens e Adultos e de ensino noturno regular ; c ) Determinação de uma Base Nacional Comum Curricular que definirá direitos e objetivos de aprendizagem do ensino médio, com \textbf{oferta} por áreas do conhecimento – I – Linguagens e suas tecnologias – II – Matemática e suas tecnologias – III – Ciências da natureza e suas tecnologias – IV – Ciências humanas e sociais aplicadas ; d ) Obrigatoriedade de estudos e práticas de educação física, arte, sociologia e filosofia – excluindo -os como componentes curriculares das áreas do conhecimento, com exceção de arte que também foi previsto como componente ; e ) O ensino da língua portuguesa e da matemática obrigatório nos três anos do ensino médio ; f ) Obrigatoriedade do Inglês como língua estrangeira, facultando a \textbf{oferta} de outras línguas estrangeiras, preferencialmente o espanhol ; g ) Composição do \textbf{Currículo} formado por Base Comum Curricular ( com \textbf{carga} \textbf{horária} de até 1.800 horas ) e \textbf{Itinerários} \textbf{Formativos} ( com \textbf{carga} \textbf{horária} de, no mínimo, 1.200 horas ) ; \textbf{Itinerários} \textbf{formativos} que deverão ser organizados por meio da \textbf{oferta} de diferentes arranjos curriculares, conforme a relevância para o contexto local e a possibilidade dos sistemas de ensino, nas quatro áreas e na formação \textbf{técnica} e \textbf{profissional} ; h ) Possibilidade de concessão de certificados intermediários de \textbf{qualificação} para o trabalho, quando a formação for estruturada e organizada em etapas com terminalidade ; i ) \textbf{Profissionais} com \textbf{notório} saber reconhecido pelos respectivos sistemas de ensino, para ministrar conteúdos de áreas afins à sua formação ou experiência \textbf{profissional} ; j ) Instituição da Política de Fomento/MEC à Implementação de Escolas de Ensino Médio em Tempo Integral.
Referente à Educação Profissional, historicamente, esta foi compreendida e operacionalizada como estratégia de cristalização da divisão social do trabalho no interior das agências de formação \textbf{profissional}, na medida em que sempre foi reservada às classes menos favorecidas, estabelecendo nítida distinção entre aqueles que produziam o saber ( ensino secundário, normal e superior ) e os que executavam tarefas manuais ( ensino \textbf{profissional} ), ( CEB, 1999, p. 77 ).
O Ensino Profissionalizante, até meados dos anos 1990, como era denominado, foi \textbf{ofertado} com base nos termos previstos nas Leis Federal nº 5.692/1971 ; nº 7.044/1982 e no Parecer nº 45/1977.
Esses documentos legais estabeleciam as \textbf{diretrizes} da educação nacional, tanto para o ensino fundamental ( primeiro grau, com oito anos de \textbf{escolaridade} obrigatória ), quanto para o ensino médio ( segundo grau, não obrigatório ).
A Lei 7.044/1982, em seu artigo 1º, definia como objetivo geral “ proporcionar ao educando a formação necessária ao desenvolvimento de suas potencialidades como elemento de autorrealização, preparação para o trabalho e para o exercício consciente da cidadania ”.
A Lei Federal 7.044/1982 estabeleceu também, dentre outras \textbf{diretrizes}, as disposições básicas sobre o \textbf{currículo}, definiu o núcleo comum obrigatório em âmbito nacional para o ensino fundamental e médio e uma parte diversificada a fim de contemplar as peculiaridades locais, a especificidade dos planos dos estabelecimentos de ensino e as diferenças individuais dos alunos.
Nesse sentido, coube aos Estados a formulação de propostas curriculares que serviram de base tanto às escolas \textbf{estaduais}, como \textbf{municipais} e particulares.
A Lei de \textbf{Diretrizes} e Bases – LDB nº 9.394/96 definiu o Ensino Médio ( 2º grau não profissionalizante ) como etapa final da Educação Básica e a Educação Profissional ( 2º grau profissionalizante ) como modalidade de ensino – com \textbf{capítulo} específico nessa Lei, regulamentado pelo Decreto Federal nº 2.208/1997.
Nesse novo contexto, a educação \textbf{profissional}, regulamentada pelo Decreto Federal nº 2.208/97, tinha, dentre outros objetivos : promover a transição entre a escola e o mundo do trabalho, capacitando jovens e adultos com conhecimentos e habilidades gerais e específicas para o exercício de atividades produtivas ; especializar, aperfeiçoar e \textbf{atualizar} o trabalho em seus conhecimentos tecnológicos.
Esse mesmo Decreto definia três níveis de \textbf{oferta} da educação \textbf{profissional} : o básico – \textbf{destinado} à \textbf{qualificação} e reprofissionalização de \textbf{trabalhadores} independente da \textbf{escolaridade} prévia ; o \textbf{técnico} – \textbf{destinado} a proporcionar \textbf{habilitação} \textbf{profissional} a alunos matriculados ou egressos do ensino médio, ministrado na forma concomitante ou pós-médio ; e o tecnológico – correspondente a \textbf{cursos} de nível superior na área tecnológica, \textbf{destinados} aos egressos do ensino médio e \textbf{técnico}.
A partir da LDB nº 9.394/1996, a \textbf{regulamentação} da educação \textbf{profissional} no Brasil passou por diversos tipos de formatos e configuração da \textbf{oferta}.
O Decreto nº 5.154/2004, que \textbf{revogou} o Decreto nº 2.208/1997, propôs a organização por \textbf{cursos} e áreas, que não guardam relação com as áreas que organizam o sistema ocupacional brasileiro, a Classificação Brasileira de Ocupações – CBO.
Este último decreto traz nova configuração e, neste, os \textbf{cursos} básicos são denominados de Formação Inicial e Continuada de \textbf{trabalhadores}.
Os \textbf{cursos} de nível \textbf{técnico} são denominados \textbf{Técnico} Nível Médio, devendo ser \textbf{ofertados} nas formas concomitante ou integrado ao Ensino Médio para os egressos do ensino fundamental e subsequente ao ensino médio ( pós-médio ), para os egressos do Ensino Médio.
Por fim, a Lei nº 11.741/08 ( BRASIL, 2008a ) disciplinou as possibilidades de \textbf{oferta} nos níveis \textbf{técnico} e tecnológico.
E a recente \textbf{reforma} do ensino médio ( Lei nº 13.415 – BRASIL, 2017 ) criou uma outra possibilidade : a formação \textbf{técnica} e \textbf{profissional} – conhecida como área 5 do Novo Ensino Médio ( o 5º Itinerário ).
Importante destacar que a expansão da Rede Federal de escolas \textbf{técnicas} no Brasil ocorreu a partir de 2007, derivada de um plano coordenado de ampliação dos investimentos na educação e teve origem com o advento do Plano de Desenvolvimento da Educação ( PDE ), que propiciou a própria expansão da Rede Federal, o Programa Brasil Profissionalizado ( Decreto nº 6.302/2007 ) e, posteriormente, o Programa Nacional de Acesso ao Ensino \textbf{Técnico} e Emprego – PRONATEC ( Lei nº 12.513/2011 ).
A ampliação das escolas \textbf{federais} de EPT, que passaram de 140 entre 1999 e 2002, para 644 unidades entre 2003 e 2016 ( MAGALHÃES ; CASTIONI, 2019, p. 2 ).
Acerca da EPT, Vargas e Carzoglio ( apud MAGALHÃES ; CASTIONI, 2019 ) afirmam a necessidade de adequação dos sistemas educativos de EPT para as necessidades dos arranjos produtivos regionais e locais.
Para isso, é necessário combinar momentos na escola e momentos no trabalho para facilitar a transição ao mundo do trabalho.
O panorama apresentado acima representa um esforço para compreensão dos avanços do Ensino Médio no país, destacando as dificuldades e desafios passados e presentes na educação, em especial a este nível de ensino.
E, particularmente, possibilita a compreensão das experiências passadas em relação ao desafio de nosso Estado em reformular e implantar a nova estrutura de \textbf{currículo} amplamente discutida pela sociedade brasileira e materializada na Base Nacional Curricular Comum.

% matched lemmas: administração, adolescente, atender, atualizar, capítulo, carga, concorrer, culminar, currículo, curso, destinar, diretriz, dirigente, escolaridade, estadual, estruturante, federal, formativo, fortalecimento, habilitação, horário, itinerário, legislação, missão, municipal, média, notório, oferta, ofertante, ofertar, política, portaria, preceito, profissional, prosseguimento, qualidade, qualificação, rede, reforma, regime, regulamentação, revogar, trabalhador, trajetória, técnico
\end{textsample}
